% Standalone paragraph. Paste into your own thesis, or \input it.
% Not wired into any chapter's numbering. Lays out the challenge each cohort poses rather than
% assigning either a settled role -- which subset of which cohort ends up in the analysis is a
% scope decision not made yet. Companion to healthy_cohort_2.tex, which argues why a clinically
% referred population cannot substitute for a population-sampled one on its own.
%
% STANDING CONSTRAINT respected here: no exact figure from the Copenhagen General Population Study
% appears in this file, or anywhere under thesis/, until data access.

\subsection{Two cohorts, an open choice}
\label{sec:two-cohorts}

Which trees end up in the analysis is not decided. Each of the two cohorts this project draws on
poses a different problem, and which subset of each to keep depends on where the thesis's later
focus lands.

ImageCAS-X is a clinically referred cohort: every scan was acquired because a clinician suspected
disease, and its descriptor table records coronary artery disease present in 388 of the 800 usable
cases and absent in the remaining 412 \cite{bransby2026imagecasx,zeng2023imagecas}. The no-disease
subset is used first: for the initial work on tortuosity and branch angle, and more generally for
the early development of the topological characterization framework.

The Copenhagen General Population Study (Herlev--{\O}sterbro, CGPS) is population-sampled rather
than referred, so it carries coronary trees with and without disease in roughly the proportions the
general population does, some of them scanned twice, and it is linked to adjudicated cardiac
outcomes \cite{fuchs2023cgps}. That mixed distribution is closer to a true population sample than a
curated healthy subset, but whether it should still be filtered down to disease-free trees for a
normal-anatomy map is a question left for later. What stands out already is the repeat-scan
subset: trees that remain healthy across both scans, and trees that were healthy at the first scan
and had become diseased by the second. 